\subsection{The Adaptive Inverse Kinematics Solver}\label{subsec:ik_impl}

The analytical core of the pipeline recovers \acrshort{mano} joint rotations from
the 21 landmarks in a single pass, with no iteration and no learned component, as
formalised in \autoref{subsec:aik}. Its implementation is split across two
modules: \texttt{armatures}, which defines the skeleton and the bookkeeping that
reconciles the two landmark conventions, and \texttt{AIK}, which performs the
solve itself.

\subsubsection{Skeleton Definition and Index Bookkeeping}\label{subsubsec:armature}

The MediaPipe and \acrshort{mano} skeletons number their joints differently, as
shown in \autoref{fig:mediapipe-joints} and \autoref{fig:mano-joints}. The
\texttt{MANOArmature} class holds the data structures that bridge them. It
describes a skeleton of 16 movable joints and 21 keypoints. A parent table,
\texttt{snap\_parent}, gives the parent of each of the 21 keypoints in MediaPipe
ordering, encoding the topology in which the wrist is the root and the five
finger bases attach directly to it, with each finger continuing as a chain. A
second list, \texttt{kinematic\_tree}, enumerates the fifteen non-base joints for
which a per-joint rotation must be solved, ordered so that every parent precedes
its children. A mapping, \texttt{id\_to\_rot}, routes each solved rotation into
its slot in \acrshort{mano}'s sixteen-entry pose vector, absorbing the difference
in finger ordering between the two conventions. Finally, \texttt{keypoints\_ext}
lists the five mesh vertices that define the fingertips; these are the indices
consumed by the extended regressor of \autoref{subsec:forward_model}.

\subsubsection{From Landmarks to Joint Rotations}\label{subsubsec:aik_impl}

The \texttt{adaptive\_IK} function takes a template skeleton and a target
skeleton, both as $21 \times 3$ arrays in MediaPipe ordering, and returns the
joint rotations that pose the template onto the target. All twist angles are
initialised to zero.

The global rotation of the wrist is recovered first, by the point-set alignment of
\autoref{subsec:aik}. The five vectors pointing from the wrist to the finger bases
are taken from both the template and the target, their cross-covariance matrix is
formed, and its singular value decomposition yields the rotation as
$R_0 = \mathbf{V}\mathbf{U}^{\top}$~(\autoref{eq:R0}). When the vectors are
near-coplanar the decomposition can return a reflection; this is detected through
the determinant and corrected by negating the sign of the last column of
$\mathbf{V}$ before recomputing $R_0$, recovering the nearest proper rotation
exactly as prescribed by Arun et al.~\cite{arun1987leastsquares}. Because the
metacarpal plate is treated as rigid, this single rotation is assigned to the
wrist and to all five finger bases.

The remaining joints are then processed along \texttt{kinematic\_tree} in
parent-before-child order. For each joint the parent's world position is first
propagated by applying the grandparent's rotation to the template bone, as in the
forward-kinematics step of \autoref{subsec:aik}. The observed target is brought
into the parent's local frame, and the swing rotation that aligns the rest-pose
bone with the observed bone is computed from the cross product and the angle
between them~(\autoref{eq:swing}). Since the twist is held at zero for every
joint, the local rotation reduces to this swing alone~(\autoref{eq:local_rot}),
and the joint's absolute rotation follows by composition with its parent's.

The function returns a stack of rotation matrices in which the first slot holds the
absolute global rotation and the remaining fifteen hold the per-joint relative
rotations, placed through \texttt{id\_to\_rot} so that the result is already in
\acrshort{mano} pose-vector order. The entire landmark-to-rotation computation is
thus a single matrix decomposition followed by one traversal of the tree, with no
iterative refinement and no learned network in the path --- the analytical
closed-form solver positioned in \autoref{subsubsec:sota_ik_analytical}.
